
预解算子的存在由谱到零点的距离控制。以下将这一条件落实为热半群积分的范数误差。

\subsection{谱测度}
\label{sec:ghs-spectral-measure}

\begin{theorem}[谱演算及 Green 算子]
\label{thm:ghs-spectral-measure}
令 $A\geq m^2>0$ 为自伴算子。存在唯一投影值测度 $E_A$，使得
\begin{equation}\label{eq:ghs-spectral-calculus}
 A=\int_{[m^2,\infty)}\lambda\,\dd E_A(\lambda),\qquad
 A^{-s}=\int_{[m^2,\infty)}\lambda^{-s}\,\dd E_A(\lambda)\quad(s>0).
\end{equation}
对于 $f,g\in\Hil$，Green 双线性型满足

$\langle A^{-1}g, f\rangle=\int\lambda^{-1}\,\dd\langle E_A(\lambda)g, f\rangle$。
\end{theorem}

若 $\Sigma$ 紧致且 $A$ 是椭圆算子的自伴实现，预解算子 紧致，谱由有限重数的特征值 $\lambda_j\to+\infty$ 构成。选定正交归一本征函数 $u_j$，才可写熟悉的求和式
\begin{equation}\label{eq:ghs-green-spectral-sum}
 G_A(x,y)=\sum_j\frac{u_j(x)\overline{u_j(y)}}{\lambda_j}.
\end{equation}
此式首先是分布核意义下的表达；一般维度不能任意将 $x=y$ 后逐项求和。非紧致流形可能具有连续谱，\cref{eq:ghs-spectral-calculus} 才是无需假设离散本征值的形式。

\begin{proposition}[热核、预解算子 与正则化]
若 $A\geq m^2>0$，则在强算子拓扑中
\begin{equation}\label{eq:ghs-heat-green-laplace}
 A^{-1}=\int_0^\infty\ee^{-sA}\,\dd s,
 \qquad A^{-z}=\frac1{\Gamma(z)}\int_0^\infty s^{z-1}\ee^{-sA}\,\dd s
\end{equation}
对 $\operatorname{Re}z>0$ 成立。若 $A$ 有零模，须先投影到 $\ker(A)^\perp$，或者引入红外质量调节。
\end{proposition}
\begin{proof}
谱演算的积分公式允许在不选择本征基的条件下证明 Green 算子的存在。对 $A\ge m^2I>0$，先在谱参数 $\lambda\ge m^2$ 上作普通积分
\[
 \int_0^T\ee^{-s\lambda}\,\dd s=\frac{1-\ee^{-T\lambda}}{\lambda},
 \qquad
 \left|\frac{\ee^{-T\lambda}}{\lambda}\right|
 \le\frac{\ee^{-Tm^2}}{m^2}.
\]
通过谱定理，这意味着
\[
 \int_0^T\ee^{-sA}\,\dd s=A^{-1}(I-\ee^{-TA}),\qquad
 \left\|A^{-1}-\int_0^T\ee^{-sA}\,\dd s\right\|
 \le \frac{\ee^{-Tm^2}}{m^2}.
\]
右边在 $T\to\infty$ 时趋于零，因而热半群积分甚至按算子范数收敛到 $A^{-1}$。这项证明说明了严格谱隙的作用：若零落在连续谱内，通常不能把 $m^2$ 直接置零并仍声称算子范数收敛；应先明确源所在的谱子空间或给出合适的红外正则化。
\end{proof}

对紧致、无边界、光滑 $(d-1)$ 维流形 $\Sigma$，定义 $N_A(\Lambda)$ 为特征值不超过 $\Lambda$ 的个数，重数计入。对于主符号为 $h^{ab}p_ap_b$ 的拉普拉斯型算子，Weyl 渐近律给出
